\documentclass[aspectratio=169]{beamer}

\usetheme{default}
\usecolortheme{default}
\setbeamertemplate{navigation symbols}{}
\setbeamertemplate{footline}[frame number]

\usepackage[T1]{fontenc}
\usepackage{amsmath}
\usepackage{booktabs}
\usepackage{array}
\usepackage{url}

\title{VBGS Variational Inference}
\subtitle{Original paper to VBOGS implementation}
\author{VBOGS}
\date{}

\begin{document}

\begin{frame}
  \titlepage
  \vspace{-0.5em}
  \begin{center}
    \small
    Goal: explain how VBGS turns observed colored points into posterior uncertainty,
    then how VBOGS attaches that uncertainty to Octree-AnyGS anchors for
    next-best-view scoring.
  \end{center}
\end{frame}

\begin{frame}{The One-Sentence Version}
  VBGS treats each observed colored point as evidence for a Bayesian Gaussian mixture.

  \vspace{0.8em}
  It infers:
  \begin{itemize}
    \item which Gaussian component likely explains each point
    \item how confident it is about each component's position, shape, and color
    \item how much posterior uncertainty remains after seeing the data
  \end{itemize}

  \vspace{0.8em}
  VBOGS runs this process separately inside each Octree-AnyGS anchor.
\end{frame}

\begin{frame}{Why Variational Inference?}
  The exact Bayesian posterior is hard:

  \[
    p(\pi, \mu, \Sigma, z \mid \text{observed points})
  \]

  because every point could belong to any component, and each component has
  unknown parameters.

  \vspace{0.8em}
  VBGS uses a tractable approximation:

  \[
    q(\pi)\,q(z)\prod_k q(\mu_k,\Sigma_k)
  \]

  Then it updates each factor in closed form using conjugate priors.
\end{frame}

\begin{frame}{Original Paper: Observed Data}
  Each observation is split into:

  \[
    s_n = \text{spatial position of point } n
  \]
  \[
    c_n = \text{color or feature vector of point } n
  \]

  \vspace{0.8em}
  Examples:

  \begin{itemize}
    \item 2D image case: \(s_n = [u, v]\), \(c_n = [r, g, b]\)
    \item 3D volume case: \(s_n = [x, y, z]\), \(c_n = [r, g, b]\)
  \end{itemize}

  \vspace{0.8em}
  The paper's model is a Bayesian mixture over these colored observations.
\end{frame}

\begin{frame}{Original Paper: Generative Model}
  For every point \(n\):

  \[
    z_n \sim \mathrm{Cat}(\pi)
  \]
  \[
    s_n \mid z_n = k \sim \mathrm{MVN}(\mu_{k,s}, \Sigma_{k,s})
  \]
  \[
    c_n \mid z_n = k \sim \mathrm{MVN}(\mu_{k,c}, \Sigma_{k,c})
  \]

  \vspace{0.8em}
  Meaning:
  \begin{itemize}
    \item choose a component \(k\)
    \item draw the point's position from component \(k\)
    \item draw the point's color from the same component \(k\)
  \end{itemize}
\end{frame}

\begin{frame}{Variable Glossary}
  \small
  \begin{tabular}{p{0.24\linewidth} p{0.66\linewidth}}
    \toprule
    Symbol & Meaning \\
    \midrule
    \(n\) & data point index \\
    \(k\) & Gaussian component index \\
    \(z_n\) & latent assignment for point \(n\) \\
    \(\pi\) & mixture weights over components \\
    \(s_n\) & spatial coordinate \\
    \(c_n\) & color or feature vector \\
    \(\mu_{k,*}\) & component mean \\
    \(\Sigma_{k,*}\) & component covariance \\
    \bottomrule
  \end{tabular}
\end{frame}

\begin{frame}{Original Paper: Bayesian Parameters}
  VBGS puts priors over the parameters:

  \[
    \pi \sim \mathrm{Dirichlet}(\alpha_0)
  \]
  \[
    (\mu_{k,s}, \Sigma_{k,s}) \sim \text{Normal-Wishart / Normal-Inverse-Wishart}
  \]
  \[
    (\mu_{k,c}, \Sigma_{k,c}) \sim \text{color posterior}
  \]

  \vspace{0.8em}
  So each Gaussian is not just a point estimate.

  \vspace{0.8em}
  It has a posterior distribution describing how confident the model is about
  that component.
\end{frame}

\begin{frame}{The Variational Posterior}
  VBGS approximates the true posterior with factorized pieces:

  \[
    q(z)\,q(\pi)\prod_k q(\mu_{k,s},\Sigma_{k,s})\,q(\mu_{k,c},\Sigma_{k,c})
  \]

  \vspace{0.8em}
  This is mean-field variational inference.

  \vspace{0.8em}
  The approximation is chosen so each update is closed form:

  \begin{itemize}
    \item \(q(z)\) from expected log probabilities
    \item \(q(\pi)\) from soft assignment counts
    \item \(q(\mu,\Sigma)\) from weighted sufficient statistics
  \end{itemize}
\end{frame}

\begin{frame}{Step 1: Soft Assignments}
  For each point and component, VBGS computes a score:

  \[
    \mathrm{score}_{nk}
    =
    \mathbb{E}_q[\log \pi_k]
    + \mathbb{E}_q[\log p(s_n \mid k)]
    + \mathbb{E}_q[\log p(c_n \mid k)]
  \]

  Then:

  \[
    r_{nk} = q(z_n = k) = \mathrm{softmax}_k(\mathrm{score}_{nk})
  \]

  \vspace{0.8em}
  \(r_{nk}\) is the responsibility: how much component \(k\) explains point \(n\).
\end{frame}

\begin{frame}{Step 2: Weighted Counts}
  Responsibilities act like fractional data assignments.

  \[
    N_k = \sum_n r_{nk}
  \]

  For mixture weights:

  \[
    \alpha_k = \alpha_{0k} + N_k
  \]

  \vspace{0.8em}
  If many points choose component \(k\), its Dirichlet posterior gets stronger.
\end{frame}

\begin{frame}{Step 3: Weighted Gaussian Stats}
  For each component, accumulate weighted sufficient statistics:

  \[
    \sum_n r_{nk}
  \]
  \[
    \sum_n r_{nk}s_n,\quad \sum_n r_{nk}s_ns_n^\top
  \]
  \[
    \sum_n r_{nk}c_n,\quad \sum_n r_{nk}c_nc_n^\top
  \]

  These update the Normal-Wishart style posterior for each component.

  \vspace{0.8em}
  Dense, consistent evidence lowers uncertainty. Sparse or inconsistent
  evidence keeps uncertainty high.
\end{frame}

\begin{frame}{Step 4: ELBO}
  The evidence lower bound is:

  \[
    \mathrm{ELBO}
    =
    \text{expected data fit}
    - \mathrm{KL}(q(\pi)\|p(\pi))
  \]
  \[
    \hspace{5.8em}
    - \mathrm{KL}(q(\text{spatial params})\|p(\text{spatial params}))
    - \mathrm{KL}(q(\text{color params})\|p(\text{color params}))
  \]

  \vspace{0.8em}
  In the paper, this is the optimization target.

  \vspace{0.8em}
  In VBOGS, it also helps decide whether an anchor needs more mixture components.
\end{frame}

\begin{frame}{What Makes This Useful}
  VBGS does not only estimate a scene.

  \vspace{0.8em}
  It estimates confidence in the scene:

  \begin{itemize}
    \item uncertain component means
    \item uncertain component covariances
    \item uncertain mixture existence or weight
    \item ambiguous point-to-component assignment
  \end{itemize}

  \vspace{0.8em}
  That posterior uncertainty is the signal VBOGS wants for next-best-view planning.
\end{frame}

\begin{frame}[fragile]{VBOGS Adaptation: Scene First}
  VBOGS starts with Octree-AnyGS:

\begin{verbatim}
posed RGB frames -> Octree-AnyGS anchors
\end{verbatim}

  Anchors are the spatial support of the rendered scene.

  \vspace{0.8em}
  Then stereo creates a colored point cloud:

\begin{verbatim}
stereo pairs -> points_world.npz
points_world = [x, y, z, r, g, b]
\end{verbatim}

  VBGS is used as a post-hoc uncertainty head.
\end{frame}

\begin{frame}[fragile]{VBOGS Coordinate Split}
  VBOGS keeps two versions of the point cloud:

\begin{verbatim}
points_world: used for anchor bucketing
points_norm:  used for VBGS fitting
\end{verbatim}

  Why?

  \begin{itemize}
    \item Octree-AnyGS anchors live in world coordinates.
    \item VBGS priors expect normalized data.
  \end{itemize}

  Mixing these frames would put points into the wrong anchors or give VBGS
  badly scaled inputs.
\end{frame}

\begin{frame}[fragile]{VBOGS Anchor Bucketing}
  Each point is assigned to every anchor cell that contains it across LOD levels:

\begin{verbatim}
cell_size(level) = voxel_size / fork^level
grid_coord = round((xyz_world - init_pos) / cell_size)
\end{verbatim}

  Important difference from a finest-only scheme:

  \begin{itemize}
    \item coarse anchors also receive points
    \item far-camera LOD anchors get realistic uncertainty
    \item well-observed regions do not look falsely uncertain at coarse levels
  \end{itemize}
\end{frame}

\begin{frame}[fragile]{VBOGS Per-Anchor Fit}
  For each anchor \(a\):

\begin{verbatim}
pts_a = normalized 6D points assigned to anchor a
\end{verbatim}

  If too few points:

\begin{verbatim}
anchor is unobserved -> U_MAX later
\end{verbatim}

  Otherwise:

\begin{verbatim}
Fit VBGS mixture on pts_a
Save posterior parameters for anchor a
\end{verbatim}

  This is Stage M4b in \path{scripts/fit_anchors.py}.
\end{frame}

\begin{frame}{VBOGS: Mapping Paper Variables}
  \small
  \begin{tabular}{p{0.37\linewidth} p{0.52\linewidth}}
    \toprule
    Paper & VBOGS implementation \\
    \midrule
    one global set of observations & points assigned to one anchor \\
    \(s_n\) & \texttt{points\_norm[:, :3]} \\
    \(c_n\) & \texttt{points\_norm[:, -3:]} \\
    \(z_n\) & local component assignment inside anchor \\
    \(\pi\) & anchor-local mixture weights \\
    \(K\) & anchor-local component count \\
    posterior entropy & anchor uncertainty scalar \(U[a]\) \\
    \bottomrule
  \end{tabular}
\end{frame}

\begin{frame}[fragile]{VBOGS: Code-Level VI Loop}
  In \path{vbgs/vbgs/model/train.py}:

\begin{verbatim}
space_logprob = expected_log_likelihood(spatial)
color_logprob = expected_log_likelihood(color)
prior_logprob = log_mean(pi)

responsibilities = softmax(space + color + prior)
\end{verbatim}

  Then sufficient stats are accumulated and conjugate updates are applied:

\begin{verbatim}
posterior params = prior params + weighted stats
\end{verbatim}
\end{frame}

\begin{frame}[fragile]{VBOGS: Batched Anchor Variant}
  The repo adds a batched path for speed.

  \vspace{0.8em}
  Instead of fitting one anchor at a time:

\begin{verbatim}
data_padded: [anchors, points, 6]
valid_mask:  [anchors, points]
\end{verbatim}

  The code transposes point samples to the front, computes responsibilities per
  anchor, masks padded rows, and updates all anchors in the group.

  \vspace{0.8em}
  This preserves VBGS math while improving throughput.
\end{frame}

\begin{frame}[fragile]{VBOGS: K Growth}
  VBOGS starts with:

\begin{verbatim}
K_INIT = pc.n_offsets
\end{verbatim}

  Then it tries larger mixtures:

\begin{verbatim}
K -> K * K_GROWTH_FACTOR
\end{verbatim}

  It accepts larger \(K\) only when mean per-point ELBO improves by enough:

\begin{verbatim}
ELBO_next - ELBO_current >= ELBO_IMPROVEMENT_TOL
\end{verbatim}

  Caveat: ELBO-as-model-selection is pragmatic, not perfect.
\end{frame}

\begin{frame}[fragile]{Saved Posterior Artifact}
  \path{anchor_posterior.npz} stores, per observed anchor:

\begin{verbatim}
alpha
spatial_mean, spatial_kappa, spatial_u, spatial_n
delta_mean,   delta_kappa,   delta_u,   delta_n
final_k
is_observed
fit_completed
\end{verbatim}

  These are posterior parameters, not rendered colors.

  \vspace{0.8em}
  The uncertainty step reads this artifact.
\end{frame}

\begin{frame}{From Posterior to Uncertainty}
  For each active component:

  \[
    \text{component\_entropy}_k
    =
    \text{Normal-Wishart entropy of spatial posterior}
    +
    \text{delta/color posterior entropy}
  \]

  Component weights come from:

  \[
    \text{weight}_k = \frac{\alpha_k}{\sum_j \alpha_j}
  \]

  Anchor uncertainty:

  \[
    U[a] = \sum_k \text{weight}_k \cdot \text{component\_entropy}_k
  \]

  Current project choice: Dirichlet entropy is saved as a diagnostic, but it is
  not added to \(U[a]\).
\end{frame}

\begin{frame}{What U Means}
  High \(U[a]\) can mean:

  \begin{itemize}
    \item sparse evidence in that anchor
    \item noisy or inconsistent spatial samples
    \item mixed geometry that needs more components
    \item uncertain color/material observations
    \item unobserved anchor assigned \path{U_MAX}
  \end{itemize}

  \vspace{0.8em}
  Low \(U[a]\) means:

  \begin{itemize}
    \item enough points
    \item consistent geometry/color
    \item confident posterior parameters
  \end{itemize}
\end{frame}

\begin{frame}[fragile]{Next-Best View Link}
  VBOGS renders uncertainty through Octree-AnyGS geometry:

\begin{verbatim}
normal render: anchors -> color image
scalar render: anchors -> uncertainty image
\end{verbatim}

  For each candidate pose:

\begin{verbatim}
score = sum(uncertainty_image) / (sum(alpha_image) + EPS)
\end{verbatim}

  The selected pose sees the highest uncertainty per visible surface area.
\end{frame}

\begin{frame}{Paper vs Implementation}
  \scriptsize
  \begin{tabular}{p{0.18\linewidth} p{0.36\linewidth} p{0.36\linewidth}}
    \toprule
    Aspect & Original VBGS paper & VBOGS repo \\
    \midrule
    Scope & a Bayesian Gaussian splatting scene/model & uncertainty head over Octree-AnyGS anchors \\
    Data & image/volume observations with position and color & stereo 6D point cloud \\
    Fit unit & mixture components for scene representation & one local mixture per anchor \\
    Output & posterior over Gaussian parameters & \path{U.npy} plus NBV scores \\
    Renderer role & VBGS can render its mixture & Octree-AnyGS renders; VBGS supplies uncertainty \\
    \bottomrule
  \end{tabular}
\end{frame}

\begin{frame}{Key Implementation Files}
  Read in this order:

  \begin{enumerate}
    \item \path{docs/Algorithm.txt}
    \item \path{docs/ALGORITHM_OVERVIEW.md}
    \item \path{docs/COLORED_POINT_CLOUD_UNCERTAINTY.md}
    \item \path{scripts/bucket_points.py}
    \item \path{scripts/fit_anchors.py}
    \item \path{vbgs/vbgs/model/train.py}
    \item \path{scripts/compute_uncertainty.py}
    \item \path{vbogs/render.py}
    \item \path{scripts/score_nbv.py}
  \end{enumerate}
\end{frame}

\begin{frame}[fragile]{Takeaways}
  VBGS is the Bayesian inference engine:

\begin{verbatim}
points -> responsibilities -> posterior parameters -> entropy
\end{verbatim}

  VBOGS is the robotics/scene wrapper:

\begin{verbatim}
Octree anchors + stereo points -> per-anchor VBGS -> U.npy -> NBV score
\end{verbatim}

  The crucial conceptual move:

\begin{verbatim}
uncertainty is not a learned image channel;
it is derived from posterior distributions over local Gaussian parameters.
\end{verbatim}
\end{frame}

\begin{frame}{References}
  \begin{itemize}
    \item Local paper copy: \path{docs/VBGS-Paper.pdf}
    \item Paper: Variational Bayes Gaussian Splatting, arXiv:2410.03592
    \item Algorithm spec: \path{docs/Algorithm.txt}
    \item Implementation plan: \path{PLAN.md}
    \item Original VBGS code: \path{vbgs/vbgs/model/train.py}
    \item VBOGS anchor fitting: \path{scripts/fit_anchors.py}
    \item VBOGS uncertainty scalar: \path{scripts/compute_uncertainty.py}
  \end{itemize}
\end{frame}

\end{document}
